% !TEX root = ../../ctfp-print.tex

\begin{quote}
  しばらく前から、プログラマーを対象とした圏論の本を書くというアイデアを
  温め続けていた。コンピュータ科学者ではなく、プログラマー向けに――
  科学者ではなくエンジニア向けに、だ。これが無謀に聞こえることは承知している
  し、正直なところ怖い。科学とエンジニアリングの間には大きな溝があることは
  否定できない。私はその両方の側で働いてきたからだ。しかし、私はいつも物事を
  説明したいという強い衝動を感じてきた。シンプルな説明の達人であるリチャード・
  ファインマンには多大な敬意を抱いている。私はファインマンではないが、最善を
  尽くすつもりだ。まずはこの前書きを公開することから始める――読者が圏論を
  学ぶ動機付けとなることを目的として――議論を始め、フィードバックを募ることを
  期待して。\footnote{
    この内容を実際の聴衆に教えている様子を、
    \href{https://goo.gl/GT2UWU}{https://goo.gl/GT2UWU} で視聴することもできる
    （またはYouTubeで「bartosz milewski category theory」と検索してほしい）。}
\end{quote}

\lettrine[lhang=0.17]{数}{段落にわたって}、この本があなたのために書かれたもの
であると納得してもらえるよう試みる。数学の最も抽象的な分野のひとつを「あり余る
余暇時間」に学ぶことへのどんな異議も、まったく根拠がないことを。

私の楽観主義はいくつかの観察に基づいている。第一に、圏論は非常に有用な
プログラミングアイデアの宝庫である。Haskellプログラマーは長い間この資源を
活用してきており、そのアイデアは徐々に他の言語にも浸透しつつあるが、この
プロセスは遅すぎる。加速させる必要がある。

第二に、数学にはさまざまな種類があり、それぞれが異なる聴衆に訴えかける。
微積分や代数に馴染めなくても、圏論を楽しめないということにはならない。
私は圏論がプログラマーの思考に特に適した数学であると主張したい。なぜなら、
圏論は個別の事物を扱うのではなく、構造を扱うからだ。プログラムを合成可能に
する種類の構造を扱う。

合成は圏論の根幹にある――圏そのものの定義の一部だ。そして私は、合成が
プログラミングの本質であると強く主張する。私たちはずっと物事を合成してきた。
偉大なエンジニアがサブルーチンのアイデアを思いつくずっと前から。かつて
構造化プログラミングの原則がプログラミングに革命をもたらした。コードの
ブロックを合成可能にしたからだ。その後にオブジェクト指向プログラミングが
登場した。これはオブジェクトの合成を中心としたものだ。関数型プログラミングは
関数や代数的データ構造の合成だけではない――並行性を合成可能にする――これは
他のプログラミングパラダイムではほぼ不可能なことだ。

第三に、私には秘密兵器がある。肉切り包丁だ。それで数学を叩き切って、
プログラマーが受け入れやすいものにする。専門的な数学者であれば、すべての
前提を正確に整理し、すべての命題を適切に限定し、すべての証明を厳密に構成
しなければならない。これが数学の論文や本を部外者にとって非常に読みにくい
ものにしている。私は物理学者の訓練を受けており、物理学では非形式的な推論を
使って驚くべき進歩を遂げてきた。数学者たちはディラックのデルタ関数を笑った。
それは偉大な物理学者P・A・M・ディラックがいくつかの微分方程式を解くために
即興で作り出したものだ。彼らは笑いをやめた。ディラックの洞察を形式化した
超関数論という全く新しい解析学の分野を発見したときに。

もちろん、手を振るような議論を使う場合は明らかに誤ったことを言うリスクがある。
そのため、この本の非形式的な議論の背後には確固たる数学理論があることを
確認するよう努める。サンダース・マクレーンの\emph{Categories for the Working
  Mathematician}の使い古しのコピーが枕元にある。

これはプログラマー\emph{のための}圏論であるから、すべての主要な概念を
コンピュータコードを使って説明する。関数型言語が一般的な命令型言語よりも
数学に近いことはご存知だろう。また、より高い抽象化能力を持つ。だから自然な
誘惑として、こう言いたくなるだろう：圏論の恩恵を受けるためにはHaskellを
学ばなければならない、と。しかしそれは圏論が関数型プログラミング以外に
応用がないことを意味するが、それは単純に間違いだ。そのため、多くのC++の例を
提供する。確かに、醜い構文を克服しなければならないし、冗長さの背景から
パターンが浮かび上がらないかもしれないし、より高い抽象化の代わりに
コピー\&ペーストを強いられるかもしれない。しかし、それがC++プログラマーの
宿命だ。

しかし、Haskellに関しては免除されるわけではない。Haskellプログラマーになる
必要はないが、C++で実装するアイデアをスケッチして文書化するための言語として
必要だ。それがまさに私がHaskellを始めたきっかけだ。その簡潔な構文と強力な
型システムが、C++のテンプレート、データ構造、アルゴリズムを理解し実装する上で
大いに役立つことを発見した。しかし、読者がすでにHaskellを知っていることを
期待することはできないので、ゆっくりと導入し、進むにつれてすべてを説明する。

経験豊富なプログラマーなら、こう自問するかもしれない：圏論や関数型の手法を
気にすることなくずっとコーディングしてきたのに、何が変わったのか、と。
新しい関数型の機能が命令型言語に次々と侵入してきていることに気づかずには
いられないだろう。オブジェクト指向プログラミングの牙城であるJavaでさえ、
ラムダを受け入れた。C++は最近、急速な速度で進化しており――数年ごとに新しい
標準が出ている――変化する世界に追いつこうとしている。これらすべての活動は、
破壊的な変化、物理学者が言うところの相転移に向けた準備だ。水を加熱し続ければ、
いずれは沸騰し始める。私たちは今、ますます熱くなる水の中で泳ぎ続けるべきか、
それとも別の選択肢を探し始めるべきかを判断しなければならないカエルの立場にある。

\begin{figure}[H]
  \centering
  \includegraphics[width=0.5\textwidth]{images/img_1299.jpg}
\end{figure}

\noindent
この大きな変化を推進する力のひとつがマルチコア革命だ。現在主流のプログラミング
パラダイムであるオブジェクト指向プログラミングは、並行性や並列性の領域では
何も得られず、むしろ危険でバグの多い設計を促進する。オブジェクト指向の基本前提
であるデータ隠蔽は、共有と変更と組み合わせると、データ競合のレシピになる。
ミューテックスとそれが保護するデータを組み合わせるアイデアは良いが、残念ながら
ロックは合成できず、ロックの隠蔽はデッドロックをより発生しやすくデバッグを
難しくする。

しかし並行性がない場合でも、ソフトウェアシステムの複雑さの増大が命令型
パラダイムのスケーラビリティの限界を試している。簡単に言えば、副作用が
手に負えなくなってきている。確かに、副作用を持つ関数は便利で書きやすいことが
多い。その効果は原則として名前やコメントでエンコードできる。SetPasswordや
WriteFileという名前の関数が何らかの状態を変更して副作用を生成していることは
明らかで、私たちはそれを扱うことに慣れている。問題が起きるのは、副作用を持つ
関数の上に副作用を持つ関数を合成し始め、さらにその上にというように積み重ねて
いくときだけだ。副作用が本質的に悪いわけではない――問題は、副作用が見えない
ところに隠れているため、より大きなスケールで管理することが不可能になること
だ。副作用はスケールしない。そして命令型プログラミングはすべて副作用について
のものだ。

ハードウェアの変化とソフトウェアの複雑さの増大が、プログラミングの基盤を
再考することを余儀なくさせている。ヨーロッパの偉大なゴシック大聖堂の建設者と
同様に、私たちは素材と構造の限界まで技術を磨いてきた。フランスには
\urlref{http://en.wikipedia.org/wiki/Beauvais_Cathedral}{未完成のゴシック様式の
  ボーヴェ大聖堂}があり、これが限界との深く人間的な闘いの証人として立っている。
かつての高さと軽さのあらゆる記録を打ち破ることを意図していたが、相次ぐ崩壊に
見舞われた。鉄棒や木製の支柱のような場当たり的な対策が崩壊を防いでいるが、
明らかに多くのことがうまくいかなかった。現代の観点からすると、現代の材料科学、
コンピュータモデリング、有限要素解析、および一般的な数学と物理学の助けなしに、
これほど多くのゴシック建築が成功裏に完成したのは奇跡だ。将来の世代が、複雑な
オペレーティングシステム、ウェブサーバー、インターネットインフラを構築する際に
私たちが見せてきたプログラミングスキルに同様の称賛を抱いてくれることを願う。
そして、率直に言って、そうなるべきだ。なぜなら私たちはこれをすべて非常に
脆弱な理論的基盤に基づいて行ってきたからだ。前進したいなら、その基盤を
修正しなければならない。

\begin{figure}
  \centering
  \includegraphics[totalheight=0.5\textheight]{images/beauvais_interior_supports.jpg}
  \caption{ボーヴェ大聖堂の崩壊を防ぐ場当たり的な対策。}
\end{figure}
